% TOP_PROBLEM: {"id": 30006474, "title": "Effective Mordell Bounds for Rational Points", "category_id": 1, "category": {"id": 1, "name": "number_theory", "display_name": "Number Theory", "description": "Properties of integers, prime numbers, Diophantine equations.", "slug": "number-theory", "order_index": 1, "created_at": "2026-07-31T15:26:25.670Z"}, "status": "open"}
% FIRST_POSTED: 2026-09-27
% !TeX program = lualatex
% ATTEMPT_STATUS: unresolved
\documentclass[11pt]{article}
\usepackage[margin=25mm]{geometry}
\usepackage{fontspec}
\setmainfont{Latin Modern Roman}
\usepackage{amsmath,amssymb,mathtools}
\usepackage{unicode-math}
\setmathfont{Latin Modern Math}
\usepackage{xurl,hyperref}
\hypersetup{colorlinks=true,urlcolor=blue}
\setcounter{secnumdepth}{0}
\setlength{\parindent}{0pt}
\setlength{\parskip}{5pt}
\setlength{\emergencystretch}{3em}
\begin{document}
\section*{\texorpdfstring{Effective Mordell Bounds for Rational Points}{Effective Mordell Bounds for Rational Points}}\label{30006474}
\subsection{Definitions and mathematical statement}
A logarithmic height on a projective curve can be fixed by an explicit embedding into projective space. For an explicitly given smooth projective curve $X/K$ of genus at least two, seek an effectively computable constant $B_X$ with
\[
 h(P)\le B_X\qquad(P\in X(K)),
\]
or a terminating algorithm returning the entire finite set $X(K)$. Inputs include effective equations and number-field data. A finiteness proof without computable stopping information is not an effective solution; nor is a search that keeps finding points without certifying that none remain.
\par\smallskip\textit{Formulation and concept references:} \href{https://arxiv.org/pdf/2503.10443}{[S1]}.

\subsection{Short English statement}
Can all rational points on every number-field curve of genus at least two be computed with a guaranteed stopping certificate?
\subsection{Research attempt}

\paragraph{Scope, process, and source check.}
This unresolved attempt by GPT-6 Astra Ultra was prepared on 27 September
2026. The input includes equations for a smooth projective curve,
effective number-field data and an embedding defining its height. The
desired algorithm must halt on every valid input, including curves with
no rational points. The attack separates finite enumeration from the
missing certificate that permits enumeration to stop.

Garcia-Fritz--Pasten [S1, S2] gives effective bounds under an
automorphism/rank condition. In the inspected v1 text
\href{https://arxiv.org/html/2503.10443v1}{[E3]}, Theorem 1.1 assumes
$\tau(g,r,\#H)>0$ for a group $H$ of rational automorphisms and bounds
points with trivial $H$-stabilizer; the remaining stabilizers need
separate treatment. This is not a general bound for all curves. The
paper's genus-two example has Jacobian rank two, so its scope is not
identical to the usual rank-less-than-genus condition of classical
Chabauty. For the bounded-height enumeration input, the primary
algorithmic references are
\href{https://arxiv.org/abs/1403.6501}{[E1]} and
\href{https://arxiv.org/abs/1111.4963}{[E2]}.

\paragraph{Lemma 97.1: a height bound is an actual stopping rule.}
First let $K=\mathbb Q$, and write the fixed embedding as
$X\hookrightarrow\mathbb P^m$. Every rational projective point has a
unique primitive integer representative $(a_0,\ldots,a_m)$ whose first
nonzero coordinate is positive. Its logarithmic height is
\[
                         h(P)=\log\max_i|a_i|.
\]
If an effective upper bound $B$ for $h$ is available, choose an integer
$M>e^B$. Enumerate the finitely many integer tuples with
$|a_i|\le M$, discard the zero tuple and those failing the primitive
and sign conditions, and test the homogeneous equations exactly.
The output is exactly $X(\mathbb Q)$. Indeed every output is a point
on $X$, and the bound puts every rational point in the enumeration.
There are at most $(2M+1)^{m+1}$ candidates, so the procedure terminates
without making an efficiency claim.

Conversely, if an algorithm returns the complete finite set, compute
an upper bound for the height of each output and take their maximum,
or zero for an empty output. Thus the two formulations are equivalent
as effective tasks over $\mathbb Q$, with the embedding fixed.

Over a number field the primitive-integer-tuple enumeration is not
valid: ideals and units enter projective coordinates. Instead use an
effective finite covering list for bounded-height points from [E1,
E2], then test the equations exactly and remove duplicates. If those
algorithms use relative multiplicative height, the normalization is
$H_K(P)=\exp([K:\mathbb Q]h(P))$ for absolute logarithmic height.
Choose a strictly larger certified cutoff and keep all possible
boundary candidates in the finite list. Since the assumed height
bound already covers every point on $X$, extra height candidates do
not compromise completeness after the equation test. Numerical
rounding at the cutoff must not discard a possible point.

\paragraph{A useful but conditional Jacobian reduction.}
Suppose a rational base point and explicit Abel--Jacobi embedding
$j:X\hookrightarrow J$ are supplied, together with a certified full
Mordell--Weil basis $G_1,\ldots,G_r$ modulo the known finite torsion
subgroup $T$. Write the canonical height pairing matrix as $Q$ and
assume that a positive rational $\lambda$ is certified to satisfy
$Q\succeq\lambda I$. Finally suppose a proved bound
$\widehat h(j(P))\le H$ is available for all $P\in X(K)$.

For any such point there are integers $n_i$ and $t\in T$ with
$j(P)=t+\sum_i n_iG_i$, so
\begin{equation}\label{r97:lattice}
 \lambda\sum_i n_i^2\le n^{\mathsf T}Qn
       =\widehat h(j(P))\le H.
\end{equation}
It is enough to enumerate $|n_i|\le\lceil\sqrt{H/\lambda}\rceil$
and all $t\in T$. For each resulting group element $R$, solve the
algebraic equations $j(P)=R$ over $K$. Because $j$ is an embedding,
the fibre has at most one geometric point; exact elimination and
number-field arithmetic test it. In rank zero only the torsion
enumeration remains. This is a finite algorithm under its stated
inputs.

The positivity certificate for $Q-\lambda I$ can use rational
intervals for the pairing entries and rigorously bounded Cholesky
pivots. Floating-point eigenvalues alone do not certify the cutoff.
Likewise a collection of independent points is not necessarily a full
Mordell--Weil basis: a missing finite-index coset can contain curve
points. A certified index and a corresponding treatment of cosets
would be needed to weaken that input.

This reduction exposes the first failure of the elementary height
route. Positive definiteness supplies a lower bound on the height in
terms of the coordinates, not an upper bound $H$ for points lying on
the curve. Replacing $H$ in \eqref{r97:lattice} by the largest height
seen so far has no justification. Computing more group generators or
improving the pairing precision does not resolve that missing upper
bound. The reduction also assumes a rational base point; it does not
silently choose one on an empty curve.

\paragraph{A fully effective infinite family with empty output.}
For every integer $g\ge2$ let $C_g$ be the smooth projective model of
\[
                        y^2=-(x^{2g+2}+1).
\]
The polynomial on the right is squarefree in characteristic zero:
a common root with its derivative would have $x=0$, which is not a
root. The degree-two map to $\mathbb P^1$ has $2g+2$ simple branch
points and is unramified over infinity. Riemann--Hurwitz gives
$2\operatorname{genus}(C_g)-2=-4+(2g+2)$, so the genus is $g$.

There is no real affine point because the right side is strictly
negative. Infinity must also be checked. In its chart put $u=1/x$
and $v=y/x^{g+1}$; then
\[
                       v^2=-(1+u^{2g+2}).
\]
At $u=0$ the two geometric points require $v^2=-1$, and neither is
real. The affine chart and this infinity chart cover the smooth
double cover, so $C_g(\mathbb R)$ is empty. Consequently
$C_g(K)=\varnothing$ for every number field $K$ with a real embedding.
The empty list and the vacuous bound $B_{C_g}=0$ are a complete
effective answer for these inputs, in any specified projective
embedding. This is a classical local-obstruction argument, not a
new general Mordell method.

The real-embedding hypothesis cannot be removed from that argument.
Over $\mathbb Q(i)$, for example, the affine points $(0,\pm i)$ and
the two infinity points are already visible. The family tests both
the number-field quantifier and the danger of inspecting only an
affine chart.

\paragraph{A nonempty arithmetic check with a declared finite range.}
For comparison consider the genus-two example in [S1],
\[
                         D:\quad y^2=x^6+x^4+x^2+1.
\]
Substitution gives the six affine points $(0,\pm1)$ and
$(\pm1,\pm2)$, and the infinity chart has two rational points with
$v=\pm1$. To search reduced rational $x=a/b$, $b>0$, clear
denominators:
\[
             (b^3y)^2=a^6+a^4b^2+a^2b^4+b^6.
\]
A rational number whose square is an integer is an integer: in a
reduced fraction its denominator would otherwise divide the
numerator. Hence an exact integer-square test on the right side is
necessary and sufficient for rational $y$ at that $x$.

The accompanying local check enumerates $|a|\le400$, $1\le b\le400$,
$\gcd(a,b)=1$, and finds just those six affine points; it separately
records the two infinity points. The paper proves completeness of its
example using additional arithmetic information. This attempt's
finite search does not reproduce that completeness proof, and the
number 400 is a testing limit, not a derived Mordell bound.

\paragraph{Why finiteness plus exhaustive search is insufficient.}
A clean computability model makes the logical error precise. For each
program $e$ on blank input define
\[
 S_e=\{t\ge0:e\text{ halts for the first time at step }t\}.
\]
Each $S_e$ has at most one element and has uniformly decidable
membership: simulate for exactly $t$ steps. Yet a uniform algorithm
returning an upper bound for all elements of $S_e$ would decide the
halting problem by simulating to that bound. Such an algorithm cannot
exist. This is not an encoding of programs into genus-two curves and
does not prove effective Mordell false. It demonstrates why abstract
finiteness of a uniformly searchable family does not itself supply a
uniform computable height bound.

For curves, a promising certificate search would dovetail point
enumeration with attempts to prove exhaustion using local obstructions,
height bounds or verified Jacobian computations. Every certificate
must be sound. To prove the resulting algorithm total one must further
show that every input eventually admits one of the searched
certificates. Soundness of familiar certificates is not completeness
of the certificate system.

\paragraph{Verification and first unresolved step.}
The rational-height enumeration has exact normalization and finite
bounds; the number-field version explicitly imports its enumeration
algorithm. The Jacobian lemma labels all extra inputs, including a
base point, a full group basis and the upper height bound. The empty
family checks smoothness, genus and both charts. The nonempty example
uses exact integer arithmetic but only within the reported box.

The first missing statement is a computable stopping certificate for
an arbitrary input curve without the extra automorphism/rank or local
obstruction assumptions. Concrete next targets are a certified upper
bound $H$ for a further family, an effective comparison with the
catalogue's embedding height, and a completeness theorem for a
collection of arithmetic certificates. A curve with many unsuccessful
finite searches is neither a counterexample nor evidence that its
rational points cannot be computed.

\paragraph{Final status.}
\textbf{UNRESOLVED.} The height/enumeration equivalence, conditional
Jacobian cutoff and empty real-obstruction family are established.
The universally terminating algorithm requested by the catalogue is
not obtained.

\subsection{Sources}
\begin{itemize}
\item[S1]Natalia Garcia-Fritz and Hector Pasten, "Effective Mordell for curves with enough automorphisms", arXiv:2503.10443 [math.NT], v1 dated 13 March 2025.
\url{https://arxiv.org/pdf/2503.10443}.
\textit{Formulation source.}
\item[S2]Effective Mordell for curves with enough automorphisms.
\url{https://arxiv.org/abs/2503.10443}.
\textit{Further reference; not a proof endorsement.}
\item[S3]Computing genus 2 curves over Q whose Jacobian has good reduction away from 2.
\url{https://warwick.ac.uk/fac/sci/maths/people/staff/visser/bmc2024_slides.pdf}.
\textit{Further reference; not a proof endorsement.}
\item[E1]D. Krumm, Computing points of bounded height in projective space over a number field, Mathematics of Computation 85 (2016), 423--447, arXiv:1403.6501.
\url{https://arxiv.org/abs/1403.6501}.
\textit{Primary source for effective projective bounded-height enumeration; consulted 27 September 2026.}
\item[E2]J. R. Doyle and D. Krumm, Computing algebraic numbers of bounded height, Mathematics of Computation 84 (2015), 2867--2891, arXiv:1111.4963.
\url{https://arxiv.org/abs/1111.4963}.
\textit{Primary algorithmic source, including numerical-boundary care; consulted 27 September 2026.}
\item[E3]N. Garcia-Fritz and H. Pasten, Effective Mordell for curves with enough automorphisms, arXiv:2503.10443v1 (13 March 2025), HTML text, Theorem 1.1 and Section 6.2.
\url{https://arxiv.org/html/2503.10443v1}.
\textit{Version-specific theorem and example check; consulted 27 September 2026.}
\end{itemize}
\end{document}
